\section{Discussão e Análise}
\label{sec:discussao}

Os testes unitários e as baterias automatizadas em contêineres comprovaram as
propriedades de correção formal definidas para o sistema:
\begin{enumerate}
    \item \textit{Segurança (Safety):} em todos os ensaios, apenas o nó de maior
    identificador ativo proclamou-se coordenador, e mensagens com
    \texttt{election\_id} desatualizado foram descartadas sem corromper o estado;
    \item \textit{Vivacidade (Liveness):} mesmo sob falhas simultâneas dos dois nós
    de maior prioridade, o cluster convergiu para um novo líder dentro do prazo,
    sem travamentos (\textit{deadlocks}) ou oscilações (\textit{livelocks});
    \item \textit{Detecção:} atrasos pontuais inferiores ao limiar de duas
    verificações não causaram eleições indevidas, como demonstrado na estabilidade
    do líder durante o teste de falha intermediária.
\end{enumerate}

\subsection{Complexidade de Mensagens e Inicialização}
A análise experimental confirma o custo de comunicação $\mathcal{O}(N^2)$ inerente
ao algoritmo Bully clássico no pior caso. Observa-se ainda um elevado desvio padrão
no volume de mensagens para $N=10$ (com medições variando entre 22 e 195 mensagens).
Essa alta dispersão é intrínseca à assimetria do protocolo em relação a qual nó
detecta a falha primeiro: se o maior nó sobrevivente ($N-1$) constata a
indisponibilidade antes dos demais, ele contata apenas o líder inativo e declara
vitória com sobrecarga mínima; por outro lado, se um nó subordinado de baixo
identificador detecta a falha, ele convoca todos os superiores, desencadeando a
propagação em cascata e a tempestade quadrática de mensagens prevista teoricamente.

Para mitigar esse impacto durante a subida conjunta do cluster, foi implementada
uma política de inicialização na qual apenas nós com $ID \ge ID_{coord}$ (ou quando
não há coordenador pré-definido) iniciam eleições ativas. Isso reduz o tráfego de
\textit{bootstrapping} para $\mathcal{O}(N)$, reservando a eleição completa para
quando falhas reais forem detectadas.

\subsection{Limitações do Ambiente}
Embora cada nó lógico execute isolado em um contêiner Docker independente, todos os processos residem na mesma máquina física. Essa configuração impõe limitações práticas fundamentais:
\begin{enumerate}
    \item \textbf{Domínio Único de Falha:} Os contêineres compartilham o mesmo hardware e o núcleo do sistema operacional host. Consequentemente, uma falha de hardware ou parada do daemon do Docker encerra simultaneamente 100\% dos nós do sistema, não representando nós com domínios de falha independentes;
    \item \textbf{Canal Virtual e Latência:} A comunicação via rede em ponte (\textit{bridge}) opera inteiramente em memória local do host, apresentando latência sub-milissegúndica e perda de pacotes quase nula. Em redes geograficamente distribuídas (WAN) com enlaces assimétricos e partições de rede (\textit{split-brain}), o algoritmo Bully clássico não garante consenso sem mecanismos de quórum majoritário, sendo recomendadas abordagens baseadas em Paxos ou Raft para tais cenários;
    \item \textbf{Contenção de CPU:} A disputa por ciclos de processamento e agendamento no host durante tempestades de mensagens de eleição pode atrasar o processamento de respostas RPC, tornando indispensável a calibração de \textit{deadlines} e limiares de tolerância a atrasos transitórios.
\end{enumerate}
